\documentclass[11pt]{article}
\usepackage[margin=1in]{geometry}
\usepackage[T1]{fontenc}
\usepackage{lmodern}
\usepackage{microtype}
\usepackage{booktabs}
\usepackage{longtable}
\usepackage{amsmath}
\usepackage{hyperref}
\hypersetup{colorlinks=true,linkcolor=blue,urlcolor=blue}
\newcommand{\pd}{\ensuremath{p_d}}
\newcommand{\hfive}{\ensuremath{h_{50}}}
\title{Research Progress Report: Stateful Recursive Reasoning on CCB-Learn\\\large D1 twelve-run confirmation}
\author{Pratham Lahoti}
\date{August 2026}
\begin{document}
\maketitle

\begin{abstract}
This report records the first controlled GPU study of recursive reasoning models
on D1 of CCB-Learn, a locally implemented, leakage-audited version of the
Complexity Ceiling Benchmark (CCB). We compare a direct Transformer, vanilla
Tiny Recursive Model (TRM), Deep Improvement Supervision TRM (DIS-TRM), and
our Stateful Transition Recursive Model (STRM), with three optimization seeds
per model. The headline result is scientifically interesting but not yet a
method claim: STRM is highly seed-sensitive. Its held-out-depth exact
accuracies are 5.17\%, 0.83\%, and 99.17\%; direct Transformer and vanilla
TRM remain at 0\% for every seed. Seed 2 must be independently audited and
reproduced before it is interpreted as a generalization result.
\end{abstract}

\section{Question, scope, and importance}
The research question is whether a small weight-shared recursive neural model
can learn a compositional transition rule at shallow depths and execute it at
substantially greater unseen depths. This is a specialised but consequential
algorithmic-generalization problem: success requires reusable state
transformation rather than depth-specific pattern matching.

The project began from the instruction to apply TRM to CCB-style tasks.
Gaussian-splatting object deletion was explicitly removed from scope. This
report is solely about iterative/recurrent reasoning, neural state-transition
execution, and compact models that trade parameter count for repeated
computation.

\section{Sources and prior work}
\begin{itemize}
\item CCB paper: \href{https://arxiv.org/abs/2606.29278}{arXiv}; task code:
\href{https://github.com/Shubh-Chapra/Complexity_Ceiling_Benchmark}{GitHub}.
\item TRM: \href{https://arxiv.org/abs/2510.04871}{paper}; code:
\href{https://github.com/SamsungSAILMontreal/TinyRecursiveModels}{repository}.
\item DIS: \href{https://arxiv.org/abs/2511.16886}{paper}; code:
\href{https://github.com/machinestein/Deep-Improvement-Supervision}{repository}.
\item Closest comparators: \href{https://arxiv.org/abs/2604.01577}{Fast--Slow
Latent Recurrence}, \href{https://arxiv.org/abs/2409.15647}{Looped
Transformers}, and \href{https://arxiv.org/abs/1807.03819}{Universal
Transformers}.
\end{itemize}

Shared recursive updates, looped computation, and persistent latent state are
not new individually. The candidate contribution is a precise stateful
transition formulation plus controlled trace-level evaluation on CCB-Learn.
No claim that STRM is the first recurrent or persistent-latent architecture is
appropriate without a systematic literature and code audit.

\section{What has been implemented}
\begin{enumerate}
\item Deterministic D1, D2, and D3 generators; exact oracle traces; CCB
serialization/parsing; and official-record compatibility checks.
\item Leakage-resistant deterministic splits, manifests, hashes, depth and
structural-holdout infrastructure, and overlap audits.
\item Metrics for final/trace/transition exactness, retention \pd, horizon
\hfive, first divergence $k^*$, TFBC, depth AUC, and confidence intervals.
\item Direct Transformer, GRU/LSTM, looped Transformer, vanilla TRM,
DIS-TRM, fast--slow recurrence, STRM, and a D3 GNN baseline.
\item CPU tests, tiny-data overfit tests, configuration/logging/checkpoint
pipelines, and atomic Kaggle persistence.
\end{enumerate}

Before model training, the local oracle exactly reproduced all 400 supplied
official records for each of D1, D2, and D3. This establishes compatibility
with those records but does not make CCB-Learn an identical reconstruction of
unreleased upstream generators or tests.

\section{Method and protocol}
\subsection{STRM}
STRM carries an explicit answer-state estimate $y$ and latent working state
$z$ across the operation sequence. For each operation it applies shared inner
updates with recall of the initial state and current operation, then emits the
next full structured state. Vanilla TRM supplies recursive refinement; STRM
makes the transformed answer state persistent and explicit while retaining a
latent workspace.

\subsection{D1 study}
Every run used width 64, four layers/loops, batch size 64, learning rate
$10^{-3}$, 10,000 steps, and 200 bootstrap resamples. Parameter counts are
248,649 for Transformer, 123,657 for TRM/DIS-TRM, and 153,033 for STRM.
Official sealed evaluation was disabled.

The same deterministic splits were built once before the model-by-seed loop:
train/validation depths 1--20; held-out test depths 25, 30, 35, 40, 45, and
50 with 100 examples each (600 total); and strong depths 60, 80, and 100 with
100 examples each (300 total). A seed therefore changes initialization and
minibatch order, not the evaluation instances.

\section{Metrics}
Final exact accuracy requires the complete final structured state to be
correct. Transition exact accuracy measures intermediate transition correctness.
Retention \pd summarizes step-to-step survival, while \hfive is the inferred
depth at 50\% success. Depth AUC summarizes performance across the held-out
depths. These are intentionally complementary: a single early state error can
make a complete long trajectory final-exact incorrect.

\section{D1 results: all twelve runs}
\begin{table}[ht]
\centering\small
\caption{Held-out-depth 25--50 results for every run. Final and transition are
percentages. Strong final denotes final exact accuracy at depths 60--100.}
\begin{tabular}{llrrrrrr}
\toprule
Model & Seed & Final & Transition & \pd & \hfive & AUC & Strong final\\
\midrule
Transformer & 0 & 0.00 & 11.19 & 0.500 & 1.00 & 0.000 & 0.00\\
Transformer & 1 & 0.00 & 11.17 & 0.500 & 1.00 & 0.000 & 0.00\\
Transformer & 2 & 0.00 & 11.26 & 0.500 & 1.00 & 0.000 & 0.00\\
\midrule
TRM & 0 & 0.00 & 9.68 & 0.500 & 1.00 & 0.000 & 0.00\\
TRM & 1 & 0.00 & 10.51 & 0.500 & 1.00 & 0.000 & 0.00\\
TRM & 2 & 0.00 & 2.30 & 0.500 & 1.00 & 0.000 & 0.00\\
\midrule
DIS-TRM & 0 & 0.33 & 27.80 & 0.843 & 4.06 & 0.004 & 0.00\\
DIS-TRM & 1 & 0.50 & 29.78 & 0.851 & 4.28 & 0.004 & 0.00\\
DIS-TRM & 2 & 0.17 & 26.97 & 0.814 & 3.37 & 0.001 & 0.00\\
\midrule
STRM & 0 & 5.17 & 46.15 & 0.916 & 7.88 & 0.045 & 0.00\\
STRM & 1 & 0.83 & 30.54 & 0.863 & 4.71 & 0.008 & 0.00\\
STRM & 2 & 99.17 & 99.69 & 1.000 & 3104.24 & 0.990 & 92.33\\
\bottomrule
\end{tabular}
\end{table}

\begin{table}[ht]
\centering\small
\caption{Three-seed aggregation. STRM mean is dominated by its extreme
seed-2 result, so the median and range are necessary.}
\begin{tabular}{lrrrrrr}
\toprule
Model & Mean final & Median final & Final range & Mean transition & Mean \pd & Mean strong\\
\midrule
Transformer & 0.00 & 0.00 & 0.00--0.00 & 11.21 & 0.500 & 0.00\\
TRM & 0.00 & 0.00 & 0.00--0.00 & 7.49 & 0.500 & 0.00\\
DIS-TRM & 0.33 & 0.33 & 0.17--0.50 & 28.19 & 0.836 & 0.00\\
STRM & 35.06 & 5.17 & 0.83--99.17 & 58.79 & 0.926 & 30.78\\
\bottomrule
\end{tabular}
\end{table}

\section{Interpretation}
\paragraph{Baseline evidence.} Direct Transformer and vanilla TRM produce no
complete held-out-depth solutions at this budget. DIS-TRM consistently improves
local transition behavior (about 28\% transition exactness and
$\pd\approx0.836$) but does not sustain complete long trajectories.

\paragraph{STRM signal.} STRM seeds 0 and 1 partially outperform those
baselines. Seed 2 is qualitatively different: 595/600 final held-out-depth
states and 277/300 final strong-depth states are correct according to its
recorded percentages. If reproduced, that indicates a successful generalizing
optimization regime under the same architecture and budget.

\paragraph{Why caution is required.} The STRM seed range is 0.83--99.17\%.
Near-zero training loss in all STRM seeds establishes training fit, not
extrapolation. The outcome may reflect a sharp optimization phase transition,
or it may expose a bug in training, evaluation, or artifact handling. It is
therefore not SOTA evidence and not yet a stable research contribution.

\section{Verification and next steps}
The persisted-artifact audit has passed. It found 12 result records, 12 matching
checkpoints, and one shared D1 manifest hash
(e5437b07e6aaa10d19b5c211e97ab0dcf7db21dfeb657c5d42eac87fbf561893).
The STRM seed-2 checkpoint is at step 10,000 with the intended width-64,
four-loop configuration. A new evaluation process loaded that checkpoint and
exactly reproduced its stored 99.17\% final exact and 99.69\% transition
accuracy at held-out depths.

This validates the stored artifact, not the stability of training. Required
next checks are a fresh-output rerun of STRM seed 2, deterministic-algorithm
settings where feasible, negative controls using permuted operations and
labels, structural holdout, and then D2/D3 studies. For D2/D3, first perform
one durable calibration seed. Then use two independent GPU-pinned workers,
each with isolated atomic records and a shared durable ledger; this improves
wall-clock time but does not reduce total GPU-quota consumption. Long matrices
should use Kaggle Version/batch execution rather than an unattended draft
session.

\section{Limitations}
This is D1 only, a single scale/configuration, and generated CCB-Learn
evaluation rather than an upstream sealed benchmark. It cannot support claims
about D2, D3, arbitrary algorithmic tasks, language-model reasoning, or
general SOTA. The local operational record is in \texttt{README.md},
\texttt{docs/BENCHMARK.md}, \texttt{docs/IMPLEMENTATION.md}, and
\texttt{docs/RESEARCH.md}.
\end{document}
